Como se mostró, los resultados son inconsistentes con el valor aceptado. Esto se debe a distintos problemas
que surgieron a lo largo del experimento. En primer lugar la calibración del campo magnético.
Algo que sucedió es que una vez calibrado el campo magnético, al volver a medir el campo otro día bajo
las mismas condiciones de corriente $I$ y resistencias de los reóstatos, el valor cambiaba. En un primer
lugar se atribuyó esto a algún material ferromagnético cercano, pero al medir el campo magnético sin corriente se vió
que el valor del campo era únicamente el terrestre, por lo que se descartó esta hipótesis. Luego se atribuyó a la temperatura,
pero los campos fueron medidos con un ventilador prendido para que no se sobrecalienten las bobinas, por lo que esto fue descartado también.
Nuestra hipótesis final es que haya cambiado el valor de resistencia de alguna de las bobinas, lo que haría que cambie el
valor de $I$ y por tanto de $B.$

Esto hizo, entre otras cosas, que no sepamos realmente qué tan homogéneo es el campo magnético ni tampoco su valor, sino que se supuso que no cambió
desde la última calibración, lo que no es necesariamente cierto. Esto explica por qué la curva de área relativa de impacto en función
de la tensión de alimentación en la Figura \ref{fig:area_vs_V} no sigue dicha forma para algunos valores de $V,$ ya que para que esto suceda debería haber un
campo constante $B.$

Por último, debido a que la fuente de corriente utilizada inicialmente no suministraba un valor de $I$ confiable,
se tuvo que usar otra fuente de alimentación, la cual era comandada analógicamente por otra fuente, de modo que se introdujo un error adicional al tener que
realizar otra calibración para relacionar el valor de la tensión de alimentación con $B.$

Todos estos fueron limitantes para el experimento planteado donde se buscó recrear el Método de Busch para la determinación de la relación e/m.
Sin embargo, a pesar de aquellos, se alcanzó un resultado del mismo orden de magnitud que el valor tabulado, $\left( 6,7 \pm 0,4 \right)\times 10^{11}\,\unit{\coulomb\per\kilogram}$ (o $\left( 3,1 \pm 0,2 \right)\times 10^{11}\,\unit{\coulomb\per\kilogram}$ en el caso del primer orden de mínimos), 
contra $1,75882000838(55)\times 10^{11}\,\unit{\coulomb\per\kilogram}$. Se pudo, entonces, conocer la relación entre la carga y la masa del electrón en su magnitud,
mas no puede aún afirmarse que el experimento, tal como aquí fue planteado, sea satisfactorio para la obtención de una relación consistente con la tabulada.